% SPDX-License-Identifier: CC-BY-4.0
% =====================================================================
% Section 5 of D. Joyner, "A hitch-hiker's guide to the metaplectic trace
% formula for SL(2)", revised for the 2-fold cover (n = 2), November 2026.
%
% Stage 4 of ../revision_log.md. Sources: ../original-1995/cctrace4.tex
% (= ../book-1997/book7e.tex, Sections 5.1-5.2) and cctrace5.tex
% (= ../book-1997/book7f.tex, Sections 5.3-5.4). Referee items II.9-II.13,
% the Section 5 part of II.6 and the Section 5 minor items are fixed here;
% see the "Stage 4 notes" in ../revision_log.md, which also map the 1995
% numbers to the numbers here.
%
% This file is \input by the master; it compiles on its own through
% preview.tex. Explicit formulas are checked by checks/sec5_n2_check.py.
% Comments beginning "% [Claude]" are editorial notes by Claude.
% Revised after the second referee report (ChatGPT, November 2026): diagram
% of the objects, Lemma 5.7 (the Weil representation on the cover of A), which
% settles b' = -b in Theorem 5.8(c); comparisons with the 1995 notes moved to
% Appendix A.
% =====================================================================

\providecommand{\GA}{\overline{G(\A)}}

\section{The spectral side}\label{sec:spectral}

This section computes the spectral side of the trace formula for $\GA$,
class by class, and assembles the trace formula. The general theory is
due to Langlands~\cite{L2} and Arthur~\cite{A3,A14} for connected
reductive groups, and to M\oe glin--Waldspurger~\cite{MW2} and
Li~\cite{LiIII} for covering groups; as in Section~\ref{sec:geometric},
we state it for the double cover of $G=\Res_{k/\Q}SL(2)$ and make the
terms explicit. Three things are new compared with the 1995 notes. The
residual spectrum is not empty: it consists of one even Weil
representation for each class in $k^\times/k^{\times2}$
(Theorem~\ref{thm:residual}). These classes are exactly the cuspidal
data $(A,\mu)$ with $\mu^w=\mu$, which the 1995 notes omit
(Lemma~\ref{lem:w-action}, Theorem~\ref{thm:Jchi}(c)). And the operator
$M_B^T$ is $2T-M^{-1}M'$ (Theorem~\ref{thm:Jchi}), whose $T$-term we check
against the geometric side (Proposition~\ref{prop:Tcoef}). The
normalizing factor of the intertwining operator, $L^S(2\lambda,\chi^2)/
L^S(2\lambda+1,\chi^2)$, is computed in Lemma~\ref{lem:intertwining}.

\paragraph{Roadmap.} The objects attached to a genuine automorphic
character $\mu$ of the cover of $A$ are related as follows.
\[
\begin{array}{ccc}
\varphi\in\mathcal H_B(\mu)
 & \xrightarrow{\ \text{sum over }B(\Q)\backslash G(\Q)\ }
 & E(x,\varphi,\lambda)\\[6pt]
 & & \Big\downarrow{\scriptstyle\ \text{constant term}}\\[6pt]
\theta_\varphi=\mathrm{Res}_{\lambda=1/2}E(x,\varphi,\lambda)\in V_b
 & \xleftarrow{\ \text{pole at }\lambda=\frac12\ (\mu^w=\mu)\ }
 & \varphi_\lambda+(M(w,\lambda)\varphi)_{-\lambda}
\end{array}
\]
For $\lambda=it$ the Eisenstein series are not square integrable; their
wave packets span the continuous part $L^2_{\chi,\mathrm{cont}}$ of the
spectrum.
Section~\ref{sec:decomp} sorts the spectrum by cuspidal data $\chi$.
Section~\ref{sec:eisenstein} constructs the Eisenstein series $E$, whose
constant term defines the intertwining operator $M(w,\lambda)$. The operator factors as a quotient of partial
$L$-functions times finitely many local operators
(Lemma~\ref{lem:intertwining}); its only pole in $\mathrm{Re}\,\lambda>0$
is at $\lambda=\frac12$ and occurs exactly for the ramified classes, and
the residues of $E$ there span the residual spectrum, which consists of
even Weil representations (Theorem~\ref{thm:residual}).
Section~\ref{sec:Jchi} truncates the kernels as on the geometric side and
computes the integral $J^T_\chi(f)$ of each truncated kernel from the
inner products of truncated Eisenstein series (the Maass--Selberg
relations): the diagonal terms give $2T-M^{-1}M'$, and for a ramified
class the cross terms give $\frac14\operatorname{tr}(M(w,0)I(\mu_0,f))$.
Section~\ref{sec:tf} equates the two sides.

\paragraph{Conventions.} $L^2_{\mathrm{gen}}=L^2(G(\Q)\backslash\GA)_
{\mathrm{gen}}$ is the space of $\varphi$ with $\varphi(\gamma x\zeta)=
\zeta\varphi(x)$, with the measure of Section~\ref{sec:global} (total
volume $1$), and $R$ is the right regular representation; test functions
$f\in\mathcal H(\GA)$ are anti-genuine and $R(f)=\int f(y)R(y)\,dy$, as in
Section~\ref{sec:kernels}. As in Section~\ref{sec:chambers} and
Remark~\ref{rem:alpha}, $\mathfrak a\cong\R$ by $X\mapsto X(\chi_1)$, with
Lebesgue measure $dX$, and we identify $\mathfrak a^*_\C$ with $\C$ so that
$\lambda(X)=\lambda X$. Then $e^{\lambda(H(x))}=e^{\lambda H(x)}$,
\[
\rho=\rho_B=1,\qquad\alpha=2/d,\qquad\alpha^\vee=d,\qquad d=[k:\Q],
\]
since $\delta_B=e^{2H}$, and $\lambda\in i\R$ is written $\lambda=it$ with
the measure $dt/2\pi$, dual to $dX$. For $k=\Q$ and $x=n(u)a(t)k$ the
point $x\cdot i$ of the upper half plane has imaginary part $y=t^2$, so
$e^{(\lambda+1)H(x)}=y^{s}$ with $s=(\lambda+1)/2$. We write $c_k=
\mathrm{vol}(A(\Q)\backslash A(\A)^1)$ as in Section~\ref{sec:coarse}, and
$X_A=A(\Q)\backslash\overline{A(\A)^1}$ and $\Pi_{\mathrm{aut}}(X_A)$ as
in Section~\ref{sec:torus-cover}; $\Pi_u(H)$ is the set of irreducible
unitary genuine representations of $H$, as there.

\subsection{The spectral decomposition and cuspidal data}\label{sec:decomp}

\begin{lemma}\label{lem:spec-decomp}
$L^2_{\mathrm{gen}}=L^2_{\mathrm{disc}}\oplus L^2_{\mathrm{cont}}$, where
$L^2_{\mathrm{disc}}$ is the closed span of the irreducible closed
invariant subspaces. Each $\pi\in\Pi_u(\GA)$ occurs in
$L^2_{\mathrm{disc}}$ with finite multiplicity $m(\pi)$, and
$L^2_{\mathrm{disc}}=L^2_{\mathrm{cusp}}\oplus L^2_{\mathrm{res}}$, where
$L^2_{\mathrm{cusp}}$ is the space of cusp forms,
\[
\int_{N(\Q)\backslash N(\A)}\varphi(nx)\,dn=0\quad\text{for almost all }x,
\]
and $L^2_{\mathrm{res}}$ is its orthogonal complement in
$L^2_{\mathrm{disc}}$.
\end{lemma}

\begin{proof}
By Lemma~\ref{lem:strong}, $G(\Q)\backslash\GA/K_f^*\cong
\widetilde\Gamma\backslash\overline{G(\R)}$ with $\widetilde\Gamma$ a
lattice, so this is Langlands' theory~\cite{L2} for
$\widetilde\Gamma\backslash\overline{G(\R)}$; in adelic form,
see~\cite{MW2}, whose setting includes finite central covers. The finiteness of multiplicities in
$L^2_{\mathrm{disc}}$ is Harish-Chandra's finiteness theorem for
automorphic forms.
\end{proof}

\begin{lemma}\label{lem:factor}
Every $\pi\in\Pi_u(\GA)$ is a restricted tensor product
$\pi\cong\bigotimes'_v\pi_v$, with $\pi_v\in\Pi_u(\overline{SL(2,k_v)})$
genuine and $\pi_v$ having a nonzero $K_v^*$-fixed vector for almost all
$v$; the restricted product is taken with respect to these vectors.
Conversely every such family defines an element of $\Pi_u(\GA)$.
\end{lemma}

\begin{proof}
A representation of $\GA$ is a representation of the restricted product
of the $\overline{SL(2,k_v)}$ (with respect to the $K_v^*$) that is
trivial on $Z_0=\{(\zeta_v)\mid\prod_v\zeta_v=1\}$
(Section~\ref{sec:global}). Flath's theorem~\cite{F} applies to the
restricted product and gives $\pi\cong\bigotimes'\pi_v$. Each $\mu_2$ at
$v$ acts by a sign $\varepsilon_v$; triviality on $Z_0$ forces all
$\varepsilon_v$ to be equal, and genuineness forces them to be the
identity. Conversely, if every $\pi_v$ is genuine, $\bigotimes'\pi_v$ is
trivial on $Z_0$ and descends.
\end{proof}

\paragraph{Cuspidal data.} The \emph{cuspidal data} are the pairs
$(G,\pi)$, with $\pi\in\Pi_u(\GA)$ occurring in $L^2_{\mathrm{cusp}}$, and
the pairs $(A,\mu)$, with $\mu\in\Pi_{\mathrm{aut}}(X_A)$ a genuine
character of $X_A$ (Theorem~\ref{thm:corr}). Among the standard
parabolic subgroups, $B$ is associate only to itself, with
$W(\mathfrak a)=\Omega(\mathfrak a,\mathfrak a)=\{1,w\}$. (The opposite
$\bar B=wBw^{-1}$ also has Levi $A$ and is associate to $B$, but it is not
standard.) For $\mu$ as above put
$\mu^w(\tilde a)=\mu(w^{-1}\tilde aw)$, where $w$ stands for its image
$(w,\kappa_\A(w))=(w,1)$ in $\GA$. Two data are \emph{equivalent} if they
are equal or of the form $(A,\mu)$, $(A,\mu^w)$. Let $\mathcal X$ be the
set of classes, $\mathcal X_G$ the classes of the $(G,\pi)$ and
$\mathcal X_A$ those of the $(A,\mu)$; for $\chi\in\mathcal X_A$ we write
$\mu\in\chi$ for the members of the set $\{\mu,\mu^w\}$, which has one or
two elements. A class $\chi\in\mathcal X_A$ is \emph{unramified} if
$\mu^w\ne\mu$, and \emph{ramified} if $\mu^w=\mu$~\cite[\S4]{A3}.

\begin{lemma}\label{lem:w-action}
\begin{enumerate}[label=(\alph*)]
\item For $t\in\A_k^\times$ and $\zeta\in\mu_2$,
  $(w,1)\,(a(t),\zeta)\,(w,1)^{-1}=(a(t^{-1}),\zeta)$ in $\GA$, and the same
  holds in each $\overline{SL(2,k_v)}$.
\item $(a(t),1)^{-1}=(a(t^{-1}),\prod_v(t_v,-1)_v)$; so $w$-conjugation is
  not inversion on $\overline{A(\A)}$, and
  $\mu^w(a(t),1)=\mu(a(t),1)^{-1}\prod_v(t_v,-1)_v$.
\item $\chi_\psi^w=\chi_\psi$ for Weil's character $\chi_\psi$
  (Lemma~\ref{lem:weil}). Hence, writing $\mu=\chi_\psi\cdot(\chi\circ p)$
  with $\chi$ a character of $A(\Q)\backslash A(\A)^1$
  (Theorem~\ref{thm:corr}), $\mu^w=\chi_\psi\cdot(\chi^{-1}\circ p)$.
\item $\mu^w=\mu$ if and only if $\chi^2=1$, and then
  $\mu=\chi_{\psi_b}$ on $\overline{A(\A)^1}$ for a unique $b\in
  k^\times/k^{\times2}$. So the ramified classes in $\mathcal X_A$ are the
  classes $\chi_b$ of $(A,\chi_{\psi_b})$, $b\in k^\times/k^{\times2}$,
  each with one member, and every other class in $\mathcal X_A$ has two
  members.
\end{enumerate}
\end{lemma}

\begin{proof}
(a) At one place, with $x(\cdot)$ as in~\eqref{eq:beta}: $x(w)=1$,
$x(a(t))=t^{-1}$, $wa(t)=\begin{pmatrix}0&-t^{-1}\\t&0\end{pmatrix}$ has
$x=t$, $wa(t)w^{-1}=a(t^{-1})$ has $x=t$, and $x(w^{-1})=-1$. So by the
conjugation formula~\eqref{eq:inverse},
$\beta(w,a(t))\beta(wa(t),w^{-1})\beta(w,w^{-1})=(t,t^2)(1,-t)(1,-1)=1$.
At non-dyadic finite $v$ the correction $\kappa_v$ is $1$ on $w$,
$w^{-1}$, $a(t)$, $a(t^{-1})$ and $wa(t)$ (the lower left entry is $0$ or
of absolute value $\ge1$), so $\sigma_v=\beta_v$ on these pairs. Taking
the product over $v$ gives the global statement.

(b) $(a(t),1)(a(t^{-1}),\varepsilon)=(1,\sigma_\A(a(t),a(t^{-1}))
\varepsilon)$ and $(t_v,t_v^{-1})_v=(t_v,t_v)_v=(t_v,-1)_v$ by~\ref{H1}
and~\ref{H2}. The formula for $\mu^w$ follows from~(a).

(c) $t^{-1}=t\cdot t^{-2}$ is in the square class of $t$, so
$\gamma_v(t^{-1})=\gamma_v(t)$ by Lemma~\ref{lem:weil}(a), and
$\chi_\psi(a(t^{-1}),\zeta)=\chi_\psi(a(t),\zeta)$. By~(a),
$\mu^w(a(t),\zeta)=\mu(a(t^{-1}),\zeta)=\chi_\psi(a(t),\zeta)\,
\chi(a(t))^{-1}$.

(d) By~(c), $\mu^w=\mu$ if and only if $\chi^2=1$. A quadratic character
of $k^\times\backslash\A_k^1$ extends uniquely to a quadratic character of
$k^\times\backslash\A_k^\times$ trivial on $\R_{>0}$ (embedded diagonally
at the archimedean places). By class field theory the quadratic
characters of $k^\times\backslash\A_k^\times$ are the
$t\mapsto\prod_v(b,t_v)_v$, $b\in k^\times$, and two values of $b$ give
the same character exactly when they differ by an element of
$k^\times\cap\prod_vk_v^{\times2}=k^{\times2}$; each such character is
trivial on $\R_{>0}$, since $(b,u)_v=1$ for $u>0$ at a real $v$. By
Lemma~\ref{lem:weil}(d), $\chi_\psi\cdot(b,\cdot)=\chi_{\psi_b}$.
\end{proof}

For $SL(2)$ itself the ramified classes are $(A,\mu)$ with $\mu^2=1$,
because there $w$ acts on $A$ by inversion~\cite{LL}. For the cover,
inversion is twisted by $(t,-1)$, and Weil's character absorbs the twist.

\begin{example}\label{ex:ramified-Q}
Let $k=\Q$. Then $k^\times/k^{\times2}$ is represented by the squarefree
integers $b$, positive and negative, and the ramified classes are the
$\chi_{\psi_b}$. If $f=f_S\otimes e^S$ as in Section~\ref{sec:local-dist},
only the classes with $\chi_{\psi_b,p}$ trivial on $A(\Z_p)$ for every
$p\notin S$ can contribute; since $(b,u)_p=1$ for all units $u$ exactly
when $\ord_p(b)$ is even, these are the $b$ whose prime factors lie in
$S$. For $S=\{\infty,2\}$ they are $b\in\{\pm1,\pm2\}$.
\end{example}

\begin{proposition}\label{prop:Lchi}
$L^2_{\mathrm{gen}}=\widehat\bigoplus_{\chi\in\mathcal X}L^2_\chi$, an
orthogonal decomposition into invariant subspaces, where:
\begin{enumerate}[label=(\alph*)]
\item for $\chi$ the class of $(G,\pi)$, $L^2_\chi$ is the
  $\pi$-isotypic subspace of $L^2_{\mathrm{cusp}}$, of multiplicity
  $m(\pi)$;
\item for $\chi\in\mathcal X_A$ with $\mu\in\chi$,
  $L^2_\chi=L^2_{\chi,\mathrm{res}}\oplus L^2_{\chi,\mathrm{cont}}$, where
  $L^2_{\chi,\mathrm{cont}}$ is the closed span of the wave packets
  $\int_\R E(x,\Phi(t),it)\,dt$ of the Eisenstein series of
  Section~\ref{sec:eisenstein}, $\Phi\in C_c^\infty(\R,\mathcal H_B(\mu))$
  with values in a finite sum of $\overline{K_0}$-types, and
  $L^2_{\chi,\mathrm{res}}=L^2_{\mathrm{res}}\cap L^2_\chi$ is spanned by
  residues of these Eisenstein series at poles with $\mathrm{Re}\,\lambda
  >0$. In particular $L^2_{\mathrm{cusp}}=\widehat\bigoplus_{\chi\in
  \mathcal X_G}L^2_\chi$.
\end{enumerate}
\end{proposition}

This is Langlands' decomposition~\cite[Ch.~7]{L2} along cuspidal
data~\cite[Ch.~II, V]{MW2}. By Lemma~\ref{lem:w-action}(d) and
Lemma~\ref{lem:intertwining}(c) below, $L^2_{\chi,\mathrm{res}}=0$ unless
$\chi$ is ramified.

\subsection{Eisenstein series and intertwining operators}\label{sec:eisenstein}

For $\mu\in\Pi_{\mathrm{aut}}(X_A)$ let $\mathcal H_B(\mu)$ be the space of
genuine functions $\varphi$ on $N(\A)A(\Q)A_B(\R)^0\backslash\GA$ that are
right $\overline{K_0}$-finite and satisfy
\begin{equation}\label{eq:HBmu}
\varphi(\tilde a x)=\mu(\tilde a)\varphi(x)\quad(\tilde a\in
\overline{A(\A)^1}),\qquad
\langle\varphi,\varphi'\rangle=\int_{X_A}\int_{\overline{K_0}}
\varphi(\tilde a\tilde k)\overline{\varphi'(\tilde a\tilde k)}\,d\tilde k\,
d\tilde a=c_k\int_{\overline{K_0}}\varphi\overline{\varphi'}.
\end{equation}
Here $A(\Q)$ and $N(\A)$ are embedded by $b\mapsto(b,1)$
(Lemma~\ref{lem:global}(b)), $\mu$ is extended to $\overline{A(\A)}$
trivially on $A_B(\R)^0$, and restriction to $\overline{K_0}$ is injective
by Lemma~\ref{lem:NAK-global}. For $\lambda\in\C$ put
$\varphi_\lambda(x)=\varphi(x)e^{(\lambda+1)H(x)}$, and let
$I(\mu_\lambda)$ be the representation on $\mathcal H_B(\mu)$ given by
$(I(\mu_\lambda,y)\varphi)_\lambda(x)=\varphi_\lambda(xy)$, and
$I(\mu_\lambda,f)=\int_{\GA}f(y)I(\mu_\lambda,y)\,dy$. It is the induced
representation $\mathrm{Ind}_{\overline{B(\A)}}^{\GA}(\mu\,e^{\lambda H})$
(normalized induction), unitary for $\lambda\in i\R$. Note that
$\varphi$ is $A_B(\R)^0$-invariant and that $\lambda$ enters only through
$\varphi_\lambda$.

For $\varphi\in\mathcal H_B(\mu)$ and $\mathrm{Re}\,\lambda>1$ put
\begin{gather}
E(x,\varphi,\lambda)=\sum_{\delta\in B(\Q)\backslash G(\Q)}
\varphi_\lambda(\delta x),\label{eq:Eis}\\
\bigl(M(w,\lambda)\varphi\bigr)_{-\lambda}(x)=\int_{N(\A)}
\varphi_\lambda(w^{-1}nx)\,dn .\label{eq:Mw}
\end{gather}
The sum is well defined because $\varphi_\lambda$ is left
$B(\Q)$-invariant (Lemma~\ref{lem:NAK-global}(b), and $\mu$ is trivial
on $A(\Q)$). By Lemma~\ref{lem:w-action}(a), $M(w,\lambda)$ maps
$\mathcal H_B(\mu)$ to $\mathcal H_B(\mu^w)$, and it intertwines
$I(\mu_\lambda)$ with $I(\mu^w_{-\lambda})$. The constant term of a
function $\phi$ on $G(\Q)\backslash\GA$ is $\phi_B(x)=\int_{N(\Q)
\backslash N(\A)}\phi(nx)\,dn$.

\begin{theorem}[Langlands~\cite{L2}; M\oe glin--Waldspurger~\cite{MW2}
for covers]\label{thm:eisenstein}
Let $\mu\in\Pi_{\mathrm{aut}}(X_A)$ and $\varphi\in\mathcal H_B(\mu)$.
\begin{enumerate}[label=(\alph*)]
\item The series~\eqref{eq:Eis} and the integral~\eqref{eq:Mw} converge
  absolutely for $\mathrm{Re}\,\lambda>1$, and $E(x,\varphi,\lambda)$
  and $M(w,\lambda)\varphi$ continue to meromorphic functions of
  $\lambda\in\C$.
\item $E_B(x,\varphi,\lambda)=\varphi_\lambda(x)+(M(w,\lambda)
  \varphi)_{-\lambda}(x)$.
\item $E(x,M(w,\lambda)\varphi,-\lambda)=E(x,\varphi,\lambda)$ and
  $M(w,-\lambda)M(w,\lambda)=1$.
\item On $i\R$, $E(x,\varphi,\lambda)$ and $M(w,\lambda)$ are holomorphic
  and $M(w,\lambda)$ is unitary. In $\mathrm{Re}\,\lambda>0$ they have at
  most finitely many poles, all simple and in $(0,1]$, and the poles of
  $E$ are those of $M(w,\lambda)$.
\item $E(x,I(\mu_\lambda,f)\varphi,\lambda)=\int_{\GA}f(y)
  E(xy,\varphi,\lambda)\,dy$, and $M(w,\lambda)I(\mu_\lambda,f)=
  I(\mu^w_{-\lambda},f)M(w,\lambda)$.
\end{enumerate}
\end{theorem}

The measure on $i\mathfrak a^*$ is $dt/2\pi$, and the number $n(A_B)$ of
chambers in $\mathfrak a$ is $2$; it produces the factor $\frac12$ in
Proposition~\ref{prop:kernels-spec}. The convergence of~\eqref{eq:Eis} for $\mathrm{Re}\,\lambda>\rho=1$ is
Godement's criterion; the integral~\eqref{eq:Mw} in fact converges for
$\mathrm{Re}\,\lambda>\frac12$, by Lemma~\ref{lem:intertwining}.

\begin{lemma}\label{lem:intertwining}
Let $\mu=\chi_\psi\cdot(\chi\circ p)\in\Pi_{\mathrm{aut}}(X_A)$. Let $S$
be a finite set of places of $k$ containing the archimedean and dyadic
places and those where $\psi_v$ has conductor $\ne\OO_v$ or $\chi_v$ is
ramified. For $v\notin S$ let $\varphi_v^0\in I(\mu_{v,\lambda})$ be the
$K_v^*$-fixed vector with $\varphi_v^0=1$ on $K_v^*$, and let $M_v(\lambda)$
be the local intertwining operator, $(M_v(\lambda)\varphi_v)(x)=
\int_{k_v}\varphi_{v,\lambda}(w^{-1}n(y)x)\,dy$ with the self-dual
measure.
\begin{enumerate}[label=(\alph*)]
\item For $\varphi=\varphi_S\otimes\bigotimes_{v\notin S}\varphi_v^0$,
  $M(w,\lambda)\varphi=(M_S(\lambda)\varphi_S)\otimes\bigotimes_{v\notin S}
  M_v(\lambda)\varphi_v^0$, with $M_S=\bigotimes_{v\in S}M_v$.
\item For $v\notin S$, with $q=q_v$ and $\chi_v(\varpi)$ the value at a
  uniformizer,
  \[
  M_v(\lambda)\varphi_v^0=\frac{1-\chi_v(\varpi)^2q^{-1-2\lambda}}
  {1-\chi_v(\varpi)^2q^{-2\lambda}}\,\varphi_v^{0,w}
  =\frac{L_v(2\lambda,\chi_v^2)}{L_v(2\lambda+1,\chi_v^2)}\,
  \varphi_v^{0,w},
  \]
  where $\varphi_v^{0,w}$ is the $K_v^*$-fixed vector of
  $I(\mu^w_{v,-\lambda})$. Hence $M(w,\lambda)=m^S(\mu,\lambda)\,M_S(\lambda)
  \otimes1$ on such $\varphi$, with
  \begin{equation}\label{eq:mS}
  m^S(\mu,\lambda)=\frac{L^S(2\lambda,\chi^2)}{L^S(2\lambda+1,\chi^2)} .
  \end{equation}
\item In $\mathrm{Re}\,\lambda>0$, $M(w,\lambda)$ is holomorphic unless
  $\chi^2=1$ ($\chi$ extended to $k^\times\backslash\A_k^\times$ trivially
  on $\R_{>0}$). If $\chi^2=1$, its only pole there is a simple pole at
  $\lambda=\frac12$, with residue $\mathrm{Res}_{s=1}\zeta_k^S(s)\,
  \bigl(2\zeta_k^S(2)\bigr)^{-1}M_S(\tfrac12)\otimes1$, which is not $0$.
\end{enumerate}
\end{lemma}

\begin{proof}
(a) The lift $(w,1)$ of $w\in G(\Q)$ is the product of the local lifts
$(w,1)_v$ ($\kappa_v(w)=1$ for all $v$), $N(\A)$ is the restricted
product of the $N(k_v)$ with $dn=\prod dn_v$, and a factorizable
$\varphi$ is a product of local functions (Lemma~\ref{lem:factor}).

(b) Since $\chi_{\psi,v}$ is trivial on $A(\OO_v)^*$
(Lemma~\ref{lem:weil}(b)), $\varphi_v^0((n,1)(a(t),\zeta)\tilde k)=\zeta
\mu_v(a(t),1)|t|^{\lambda+1}$ for $\tilde k\in K_v^*$ is well defined. If
$|y|\le1$, then $w^{-1}n(y)\in K_v$, its lift lies in $K_v^*$ because
$\sigma_v\equiv1$ on $K_v\times K_v$, and the integrand is $1$. If
$|y|=q^m$ with $m\ge1$, write $y=\varpi^{-m}u$, $u\in\OO^\times$; then
\[
w^{-1}n(y)=n(-y^{-1})\,a(-y^{-1})\,\begin{pmatrix}1&0\\y^{-1}&1
\end{pmatrix},
\]
the last factor in $K_v$. Computing the cocycle of this decomposition, the
integrand is $\chi_v(\varpi)^m|y|^{-\lambda-1}$ times $1$ if $m$ is even,
and times $c\,(u,\varpi)_v$ with a constant $c$ of absolute value $1$ if
$m$ is odd. (The odd case uses $\gamma_v(\varpi u)=(u,\varpi)_v
\gamma_v(\varpi)$, Lemma~\ref{lem:weil}(a)--(b); both cases are checked in
\texttt{checks/sec5\_n2\_check.py} at $p=3,5,7$.) The shell
$\{|y|=q^m\}$ has volume $q^m(1-q^{-1})$, and $u\mapsto(u,\varpi)_v$ is a
nontrivial character of $\OO^\times$, so the odd shells contribute $0$
and
\[
M_v(\lambda)\varphi_v^0(1)=1+(1-q^{-1})\sum_{j\ge1}\bigl(\chi_v(\varpi)^2
q^{-2\lambda}\bigr)^j=\frac{1-\chi_v(\varpi)^2q^{-1-2\lambda}}
{1-\chi_v(\varpi)^2q^{-2\lambda}} .
\]
$M_v(\lambda)\varphi_v^0$ is $K_v^*$-fixed, so it is this multiple of
$\varphi_v^{0,w}$.

(c) The local integrals at $v\in S$ converge absolutely for
$\mathrm{Re}\,\lambda>0$, since the integrand is $O(|y|^{-\mathrm{Re}\,
\lambda-1})$, so $M_S(\lambda)$ is holomorphic there. If $\chi^2\ne1$,
$L^S(s,\chi^2)$ is entire (the removed Euler factors have no poles in
$\mathrm{Re}\,s>0$) and $L^S(s,\chi^2)\ne0$ for $\mathrm{Re}\,s\ge1$.
If $\chi^2=1$, $m^S=\zeta_k^S(2\lambda)/\zeta_k^S(2\lambda+1)$, whose only
pole in $\mathrm{Re}\,\lambda>0$ is at $2\lambda=1$. The operator
$M_S(\frac12)$ is not $0$: if $\varphi_S$ is supported in
$\overline{B(\Q_S)}\,w^{-1}N(U)$ for a small compact neighborhood $U$ of
$0$, then $(M_S(\frac12)\varphi_S)_{-1/2}(1)=\int_U\varphi_{S,1/2}
(w^{-1}n(y))\,dy$, and on the big cell $\varphi_S$ may be any smooth
compactly supported function of $y$.
\end{proof}

For $SL(2)$ itself the same computation gives
$L_v(\lambda,\chi_v)/L_v(\lambda+1,\chi_v)$, with the pole at
$\lambda=1=\rho$ for $\chi=1$ (the script checks this too). On the double
cover the variable is doubled and the character squared; this is the
Gindikin--Karpelevich formula for the double cover~\cite{KP}; for
general covering groups see~\cite{McN,Gao}. The factor~\eqref{eq:mS} is
what produces the logarithmic derivative $m'/m$ in Section~6.

\paragraph{The Weil representation.} The residual spectrum consists of
Weil representations, and to say which ones we fix a normalization. Let
$v$ be a place of $k$ and $\psi'$ a nontrivial character of $k_v$. The
\emph{Weil representation} $\omega_{\psi'}$ attached to $\psi'$ and the
quadratic form $x^2$~\cite{We64,Ge2} is a genuine representation of
$\overline{SL(2,k_v)}$ on the Schwartz--Bruhat space $\mathcal S(k_v)$ in
which
\begin{equation}\label{eq:weil-rep}
\begin{gathered}
\omega_{\psi'}(n(u),1)\phi(x)=\psi'(ux^2)\phi(x),\qquad
\omega_{\psi'}(w,1)\phi(x)=c\,(\mathcal F\phi)(-x),\\
\mathcal F\phi(x)=\int_{k_v}\phi(y)\psi'(2xy)\,dy,
\end{gathered}
\end{equation}
with a constant $c$ of absolute value $1$ and the measure that is
self-dual for $\psi'(2xy)$. (In Kubota's model $c=\gamma'(1)^{-1}$, with
$\gamma'$ as in Lemma~\ref{lem:weil-rep}; see Lemma~\ref{lem:weil-kernel}.)
It is the sum of its restrictions
$r^+_{\psi'}$ and $r^-_{\psi'}$ to the even and the odd functions, which
are irreducible. In some references $n(u)$ acts by $\psi'(\frac12ux^2)$;
that representation is $\omega_{\psi'_{1/2}}$ in the present notation.

\begin{lemma}\label{lem:weil-rep}
Let $\omega=\omega_{\psi'}$ be as in~\eqref{eq:weil-rep}, and let
$\gamma'(a)$ be the Weil index of $x\mapsto\psi'(ax^2)$, as
in~\eqref{eq:weil-index}. In the model of Section~\ref{sec:local}, for
$t\in k_v^\times$ and $\zeta\in\mu_2$,
\[
\omega(a(t),\zeta)\phi(x)=\zeta\,\frac{\gamma'(1)}{\gamma'(t)}\,
|t|_v^{1/2}\phi(tx)=\zeta\,\chi_{\psi'_{-1}}(a(t),1)\,|t|_v^{1/2}\phi(tx),
\]
where $\psi'_{-1}(x)=\psi'(-x)$. Consequently $\phi\mapsto\bigl(x\mapsto
(\omega(x)\phi)(0)\bigr)$ maps $r^+_{\psi'}$ into the induced
representation $I(\chi_{\psi'_{-1},v},-\frac12)$: the function
$F(x)=(\omega(x)\phi)(0)$ satisfies $F((n,1)\tilde ax)=\chi_{\psi'_{-1}}
(\tilde a)e^{H_v(\tilde a)/2}F(x)$.
\end{lemma}

\begin{proof}
Let $\bar n(u)=w\,n(-u)\,w^{-1}$, the lower triangular unipotent matrix
with entry $u$. By Lemma~\ref{lem:local}(h),
$(\bar n(u),1)=(w,1)(n(-u),1)(w,1)^{-1}$, so by~\eqref{eq:weil-rep}
$\omega(\bar n(u),1)=\mathcal F m(-u)\mathcal F^{-1}$, where $m(u)$ is
multiplication by $\psi'(ux^2)$; the constant $c$ and the reflection
$x\mapsto-x$ of~\eqref{eq:weil-rep} cancel, because the reflection
commutes with $m(u)$. Completing the square and using~\eqref{eq:weil-index} for
$\psi'$ (with $\gamma'(1/t)=\gamma'(t)$; the measure self-dual for
$\psi'(2xy)$ is $|2|_v^{1/2}$ times the one in~\eqref{eq:weil-index}),
\[
\bigl(\mathcal Fm(1/t)\mathcal F^{-1}\phi\bigr)(x)=\gamma'(t)\,|t|_v^{1/2}
\int_{k_v}\phi(z)\,\psi'\bigl(-t(x-z)^2\bigr)\,dz .
\]
Let $w(t)=n(t)\bar n(-t^{-1})n(t)=\begin{pmatrix}0&t\\-t^{-1}&0
\end{pmatrix}$. Multiplying by $m(t)$ on both sides, the terms
$\psi'(tx^2)$ and $\psi'(tz^2)$ cancel against the square, and
\[
\omega(n(t),1)\,\omega(\bar n(-t^{-1}),1)\,\omega(n(t),1)\phi(x)
=\gamma'(t)\,|t|_v^{1/2}\,(\mathcal F\phi)(tx).
\]
Now compute in the cover with~\eqref{eq:beta}. The values of $x(\cdot)$
on $n(t)$, $\bar n(-t^{-1})$, $n(t)\bar n(-t^{-1})$ and $w(t)$ are $1$,
$-t^{-1}$, $-t^{-1}$, $-t^{-1}$, so both cocycles in the triple product are
$1$ and $(n(t),1)(\bar n(-t^{-1}),1)(n(t),1)=(w(t),1)$. Next
$(w(1),1)^{-1}=(w(1)^{-1},1)$, and $\beta(w(t),w(1)^{-1})=
\bigl(\frac{t^{-1}}{-t^{-1}},t^{-1}\bigr)_v=(-1,t)_v$, because
$w(t)w(1)^{-1}=a(t)$ has $x=t^{-1}$. Hence $(w(t),1)(w(1),1)^{-1}=
(a(t),(-1,t)_v)$ and
\[
\omega\bigl(a(t),(-1,t)_v\bigr)\phi(x)=\frac{\gamma'(t)}{\gamma'(1)}\,
|t|_v^{1/2}\phi(tx).
\]
By Lemma~\ref{lem:weil}(a) with $a=b=t$, $(\gamma'(t)/\gamma'(1))^2=
(t,t)_v=(t,-1)_v$. So $(-1,t)_v\gamma'(t)/\gamma'(1)=(\gamma'(t)/
\gamma'(1))^3=\gamma'(1)/\gamma'(t)$, which is the first formula. The
second is Lemma~\ref{lem:weil}(d) with $b=-1$:
$\chi_{\psi'_{-1}}=\chi_{\psi'}\cdot(-1,\cdot)_v=\chi_{\psi'}^{-1}$. For the
last statement, $(\omega((n(u),1)x)\phi)(0)=\psi'(0)(\omega(x)\phi)(0)$,
and $e^{(-\frac12+1)H_v(\tilde a)}=|t|_v^{1/2}$.
\end{proof}

For example, let $v$ be real and $\psi'(x)=e^{2\pi iax}$. The operator
$d\omega\begin{pmatrix}0&1\\-1&0\end{pmatrix}=i\bigl(2\pi ax^2-
\frac1{8\pi a}\frac{d^2}{dx^2}\bigr)$ generates the action of the
subgroup $\tilde k_\theta$ of Lemma~\ref{lem:explicit-r}(b); for $a>0$ its
eigenvalues on even functions are $i(\frac12+2j)$, $j\ge0$, with ground
state $e^{-2\pi ax^2}$, and for $a<0$ they are $-i(\frac12+2j)$. So
$r^+_{\psi'}$ has lowest weight $\frac12$ when $a>0$ and highest weight
$-\frac12$ when $a<0$. This agrees with the lemma:
$\chi_{\psi'_{-1}}(a(-1),1)=i$ for $a>0$, which is the condition
$m_0=\frac12$ of Lemma~\ref{lem:explicit-r}(b).

Globally, for a nontrivial character $\psi'$ of $\A_k/k$ let
$\omega_{\psi'}=\bigotimes_v\omega_{\psi'_v}$ on $\mathcal S(\A_k)$, a
genuine representation of $\GA$, and for $\phi\in\mathcal S(\A_k)$ let
\begin{equation}\label{eq:theta}
\theta^{\psi'}_\phi(x)=\sum_{\xi\in k}\bigl(\omega_{\psi'}(x)\phi\bigr)(\xi).
\end{equation}
By Weil~\cite{We64}, $\theta^{\psi'}_\phi$ is left invariant under
$G(\Q)$. By~\eqref{eq:weil-rep} its constant term is
$x\mapsto(\omega_{\psi'}(x)\phi)(0)$, because $\int_{k\backslash\A_k}
\psi'(u\xi^2)\,du=0$ for $\xi\ne0$.

\begin{theorem}\label{thm:residual}
\begin{enumerate}[label=(\alph*)]
\item $L^2_{\chi,\mathrm{res}}=0$ for every $\chi\in\mathcal X_A$ that is
  unramified.
\item For $b\in k^\times$, $V_b:=L^2_{\chi_b,\mathrm{res}}$ is the closed
  span of the residues $\theta_\varphi(x)=\mathrm{Res}_{\lambda=1/2}
  E(x,\varphi,\lambda)$, $\varphi\in\mathcal H_B(\chi_{\psi_b})$; it is
  not $0$, and $\theta_{\varphi,B}(x)=(\mathrm{Res}_{\lambda=1/2}
  M(w,\lambda)\varphi)_{-1/2}(x)$, of order $e^{H(x)/2}$.
\item $L^2_{\mathrm{res}}=\widehat\bigoplus_{b\in k^\times/k^{\times2}}V_b$,
  and each $V_b$ is irreducible. With the
  normalization~\eqref{eq:weil-rep},
  \[
  V_b\cong r^+_{\psi_{-b}}=\bigotimes\nolimits'_v r^+_{\psi_{-b,v}},
  \]
  the even Weil representation attached to $\psi_{-b}(x)=\psi(-bx)$, and
  the elements of $V_b$ are the theta functions
  $\theta^{\psi_{-b}}_\phi$ of~\eqref{eq:theta}, $\phi$ an even
  Schwartz--Bruhat function on $\A_k$.
\end{enumerate}
\end{theorem}

\begin{proof}
(a) and the first part of (b): by Proposition~\ref{prop:Lchi}(b) and
Theorem~\ref{thm:eisenstein}(d), $L^2_{\chi,\mathrm{res}}$ is spanned by
residues of $E$ at poles of $M(w,\lambda)$ with $\mathrm{Re}\,\lambda>0$,
and by Lemma~\ref{lem:intertwining}(c) there are none unless
$\chi^2=1$, i.e.\ $\mu=\chi_{\psi_b}$ (Lemma~\ref{lem:w-action}(d)), and
then only $\lambda=\frac12$. The constant term of the residue is the
residue of Theorem~\ref{thm:eisenstein}(b), in which $\varphi_\lambda$ is
holomorphic; it is nonzero by Lemma~\ref{lem:intertwining}(c), and it is
$O(e^{(-\frac12+1)H(x)})$. Since $\int_T^\infty e^{X}e^{-2X}\,dX<\infty$
for the measure $\delta_B^{-1}=e^{-2X}$, $\theta_\varphi$ is square
integrable (its difference with the constant term is rapidly decreasing
on a Siegel set).

(c) By Gelbart~\cite{Ge2}, Gelbart and
Piatetski-Shapiro~\cite{GPS80} and Waldspurger~\cite{Wa80},
$L^2_{\mathrm{res}}$ is the orthogonal sum, over the classes $b'\in
k^\times/k^{\times2}$, of the spaces $\Theta_{b'}$ of theta functions
$\theta^{\psi_{b'}}_\phi$, $\phi$ even, and $\Theta_{b'}$ is irreducible
and isomorphic to $r^+_{\psi_{b'}}$; the identification of the residues
with theta functions is the Siegel--Weil formula in this case. It remains
to match $b'$ with $b$. The constant term of an element of $V_b$
transforms on the left under $\overline{A(\A)}$ by $\chi_{\psi_b}e^{H/2}$,
by~(b) and $\chi_{\psi_b}^w=\chi_{\psi_b}$. The constant term of
$\theta^{\psi_{b'}}_\phi$ is $x\mapsto(\omega_{\psi_{b'}}(x)\phi)(0)$, which
by Lemma~\ref{lem:weil-rep}, applied at each place (on $A$ and $N$ the
cocycles $\sigma_v$ and $\beta_v$ agree), transforms by
$\chi_{\psi_{-b'}}e^{H/2}$. So $\Theta_{b'}\subset L^2_{\chi_{-b'}}$, and
therefore $\Theta_{b'}=V_{-b'}$, by the uniqueness in
Lemma~\ref{lem:w-action}(d).
\end{proof}

\begin{example}\label{ex:theta}
Let $k=\Q$, $b=1$, and $\psi$ as in Example~\ref{ex:weil-Q}, so
$\psi_{-1,\infty}(x)=e^{2\pi ix}$. The residual representation
$V_1\cong r^+_{\psi_{-1}}$ contains the theta function of
$\phi=e^{-2\pi\xi^2}\otimes\mathbf 1_{\hat\Z}$: at $x=n(u)a(y^{1/2})$ in
the real component it is $y^{1/4}\theta(u+iy)$ with the classical
$\theta(z)=\sum_{n\in\Z}e^{2\pi in^2z}$, of weight
$\frac12$ on $\Gamma_0(4)$, in agreement with the lowest weight
$\frac12$ found above. The space $V_{-1}$ contains its complex conjugate.
The multiplier of $\theta$ involves the Jacobi symbol of
Example~\ref{ex:jacobi}. In the classical variable $s=(\lambda+1)/2$ the
pole is at $s=\frac34$, and the constant term of $\theta$ is a multiple
of $y^{1/4}$. For $SL(2)$ the only residual representation is the trivial
one, the residue of $E$ at $\lambda=1$ ($s=1$), whose constant term is a
constant. With $S=\{\infty,2\}$ and $\mu=\chi_\psi$,
$m^S(\mu,\lambda)=(1-2^{-2\lambda})\zeta(2\lambda)/\bigl((1-2^{-1-2\lambda})
\zeta(2\lambda+1)\bigr)$, with residue $2/\pi^2$ at $\lambda=\frac12$.
\end{example}

The residual term is the analog for the cover of the
term~\cite[(5.5)]{LL} for $SL(2)$. The odd Weil representations, the
products $\bigotimes'_vr^{\varepsilon_v}_{\psi_{b,v}}$ with some
$\varepsilon_v=-$, are not residual: those that are automorphic occur in
$L^2_{\mathrm{cusp}}$~\cite{GPS80,Wa80}, and they are counted among the
$(G,\pi)$.

\subsection{Kernels, truncation and the terms \texorpdfstring{$J_\chi$}{J\_chi}}\label{sec:Jchi}

For $\chi\in\mathcal X$ let $K_{G,\chi}(x,y)$ be the kernel of $R(f)$ on
$L^2_\chi$, and $K_{B,\chi}(x,y)$ the kernel of $R_B(f)$
(Lemma~\ref{lem:kernels}) on the $\chi$-part of
$L^2(N(\A)A(\Q)\backslash\GA)_{\mathrm{gen}}$, which is $0$ for
$\chi\in\mathcal X_G$ and the closed span of the $\mathcal H_B(\mu)\,
e^{(\lambda+1)H}$, $\mu\in\chi$, $\lambda\in i\R$, for $\chi\in\mathcal X_A$.
For $\mu\in\Pi_{\mathrm{aut}}(X_A)$ let $\mathcal B(\mu)$ be an orthonormal
basis of $\mathcal H_B(\mu)$, for~\eqref{eq:HBmu}, consisting of
$\overline{K_0}$-finite vectors.

\begin{proposition}\label{prop:kernels-spec}
\begin{enumerate}[label=(\alph*)]
\item If $\chi$ is the class of $(G,\pi)$, then $K_{B,\chi}=0$ and
  $K_{G,\chi}(x,y)=\sum_\phi(R(f)\phi)(x)\overline{\phi(y)}$, $\phi$
  running over an orthonormal basis of $L^2_\chi$ of
  $\overline{K_0}$-finite vectors.
\item If $\chi\in\mathcal X_A$, then
  \begin{align}
  K_{B,\chi}(x,y)&=\sum_{\mu\in\chi}\frac1{2\pi}\int_\R\sum_{\varphi\in
  \mathcal B(\mu)}\bigl(I(\mu_{it},f)\varphi\bigr)_{it}(x)\,
  \overline{\varphi_{it}(y)}\,dt,\label{eq:KBchi}\\
  K_{G,\chi}(x,y)&=K_{\chi,\mathrm{res}}(x,y)+\frac12\sum_{\mu\in\chi}
  \frac1{2\pi}\int_\R\sum_{\varphi\in\mathcal B(\mu)}
  E\bigl(x,I(\mu_{it},f)\varphi,it\bigr)\,\overline{E(y,\varphi,it)}\,dt,
  \label{eq:KGchi}
  \end{align}
  where $K_{\chi,\mathrm{res}}(x,y)=\sum_\theta(R(f)\theta)(x)
  \overline{\theta(y)}$ over an orthonormal basis of
  $L^2_{\chi,\mathrm{res}}$; it is $0$ unless $\chi=\chi_b$. If
  $\mu^w\ne\mu$, the two terms of the sum over $\mu\in\chi$
  in~\eqref{eq:KGchi} are equal.
\item $\sum_{\chi\in\mathcal X}K_{P,\chi}(x,y)=K_P(x,y)=
  \sum_{\mathcal O\in\mathcal C}K_{P,\mathcal O}(x,y)$ for $P\in\{B,G\}$, the
  sum over $\chi$ converging absolutely and uniformly on compact sets.
\end{enumerate}
\end{proposition}

\begin{proof}
(a) is the spectral theorem on the discrete space $L^2_\chi$, and
$L^2_\chi$ has no constant terms. For~\eqref{eq:KBchi}, write $x=n\,
\tilde a^1e^Xk$ with $\tilde a^1\in\overline{A(\A)^1}$; for a function
$F$ on $N(\A)A(\Q)\backslash\GA$, $e^{-X}F$ is a function on
$X_A\times\R\times\overline{K_0}$ with $\int|F|^2dx=\int|e^{-X}F|^2\,
d\tilde a^1\,dX\,dk$, because $dx=\delta_B^{-1}dn\,da\,dk$. Fourier
analysis on $X_A$ (Theorem~\ref{thm:corr}(a); only genuine characters
occur) and on $\R$ (with $dX$ and $dt/2\pi$) writes $F$ as
$\sum_\mu\frac1{2\pi}\int\varphi_{\mu,t,it}\,dt$ with $\varphi_{\mu,t}\in
\mathcal H_B(\mu)$ and $\|F\|^2=\sum_\mu\frac1{2\pi}\int
\|\varphi_{\mu,t}\|^2dt$, the norm being~\eqref{eq:HBmu}. So
$R_B(f)$ is the direct integral of the $I(\mu_{it},f)$, which
gives~\eqref{eq:KBchi}. Formula~\eqref{eq:KGchi} is Langlands' spectral
decomposition of $L^2_{\chi,\mathrm{cont}}$ by wave packets of
Eisenstein series~\cite[Ch.~7]{L2}, \cite[Ch.~V]{MW2}, written as
in~\cite[\S4]{A3}; the factor $\frac12=n(A_B)^{-1}$ is the reciprocal of
the number of chambers $\pm\mathfrak a^+$ in $\mathfrak a$~\cite[p.~927]{A2}.
(Proposition~\ref{prop:Tcoef} checks this factor against the geometric
side.) The equality of the two terms is Theorem~\ref{thm:eisenstein}(c),(e)
and the unitarity of $M(w,it)$. (c) Both sides are kernels of the same
operator; the absolute convergence is~\cite[\S3]{A3}, \cite{MW2},
\cite{LiIII}.
\end{proof}

\paragraph{Truncation.} For a continuous function $\phi$ on
$G(\Q)\backslash\GA$ and $T\ge0$ put
\begin{equation}\label{eq:trunc}
\Lambda^T\phi(x)=\phi(x)-\sum_{\delta\in B(\Q)\backslash G(\Q)}
\hat\tau_B\bigl(H(\delta x)-T\bigr)\,\phi_B(\delta x).
\end{equation}
By Lemma~\ref{lem:height-ineq}(b) the sum has at most one nonzero term.

\begin{lemma}\label{lem:truncation}
Let $T\ge0$.
\begin{enumerate}[label=(\alph*)]
\item $(\Lambda^T\phi)_B(x)=0$ if $H(x)>T$.
\item $\Lambda^T\Lambda^T=\Lambda^T$, and $\langle\Lambda^T\phi,\phi'\rangle
  =\langle\phi,\Lambda^T\phi'\rangle$ whenever both sides converge
  absolutely.
\item If $\phi$ is an automorphic form, for example $E(\cdot,\varphi,
  \lambda)$ at a regular point $\lambda$, then $\Lambda^T\phi$ is rapidly
  decreasing on $G(\Q)\backslash\GA$. It is not compactly supported in
  general.
\end{enumerate}
\end{lemma}

\begin{proof}
(a) Integrating~\eqref{eq:trunc} over $n\in N(\Q)\backslash N(\A)$, the
term $\delta=1$ gives $\hat\tau_B(H(x)-T)\phi_B(x)=\phi_B(x)$. For
$\delta\ne1$ and every $n$, $H(\delta nx)\le-H(nx)=-H(x)<-T\le T$ by
Lemma~\ref{lem:height-ineq}(b), so these terms vanish. (b) Idempotence
follows from~(a), since $\Lambda^T\phi$ then has no constant term where
$\hat\tau_B(H(x)-T)\ne0$. For self-adjointness, unfold:
$\langle\phi-\Lambda^T\phi,\phi'\rangle=\int_{B(\Q)\backslash\GA}
\hat\tau_B(H(x)-T)\phi_B(x)\overline{\phi'(x)}\,dx=\int_{N(\A)B(\Q)
\backslash\GA}\hat\tau_B(H(x)-T)\phi_B(x)\overline{\phi'_B(x)}\,dx$, because
$\hat\tau_B(H(x)-T)\phi_B(x)$ is left $N(\A)$-invariant; this is symmetric.
(c) is in~\cite[\S1]{A3}; the cover splits over $N(\A)$
(Lemma~\ref{lem:global}(b)), so the proof applies verbatim.
\end{proof}

Define, for $\chi\in\mathcal X$ and $T\ge0$,
\begin{equation}\label{eq:kchi}
k^T_\chi(x,f)=K_{G,\chi}(x,x)-\sum_{\delta\in B(\Q)\backslash G(\Q)}
K_{B,\chi}(\delta x,\delta x)\,\hat\tau_B\bigl(H(\delta x)-T\bigr).
\end{equation}

\begin{lemma}\label{lem:kchi}
For $T\ge0$, $k^T_\chi(x,f)=F^G(x,T)K_{G,\chi}(x,x)+\sum_{\delta}
\tau_B(H(\delta x)-T)\bigl(K_{G,\chi}-K_{B,\chi}\bigr)(\delta x,\delta x)$,
and $\sum_\chi k^T_\chi(x,f)=k^T(x,f)$.
\end{lemma}

\begin{proof}
As for Lemma~\ref{lem:partition}, whose first identity is used; the sum
over $\chi$ is Proposition~\ref{prop:kernels-spec}(c) and~\eqref{eq:kT}.
\end{proof}

\begin{theorem}[Arthur~{\cite[\S\S2--4]{A3}}; Li~\cite{LiIII} for
covers]\label{thm:Jchi-poly}
Let $f\in\mathcal H(\GA)$.
\begin{enumerate}[label=(\alph*)]
\item For $T$ sufficiently large, $\sum_{\chi\in\mathcal X}\int_{G(\Q)
  \backslash\GA}|k^T_\chi(x,f)|\,dx<\infty$.
\item $J^T_\chi(f):=\int_{G(\Q)\backslash\GA}k^T_\chi(x,f)\,dx$ is, for
  such $T$, a polynomial in $T$, and $\sum_\chi J^T_\chi(f)=J^T(f)$, the
  geometric side of Theorem~\ref{thm:coarse}.
\item $J^T_\chi(f)-\int_{G(\Q)\backslash\GA}\Lambda^T_2K_{G,\chi}(x,x)\,dx
  \to0$ as $T\to\infty$, where $\Lambda^T_2$ is $\Lambda^T$ applied to the
  second variable.
\end{enumerate}
We put $J_\chi(f)=J^0_\chi(f)$, the value of the polynomial at $T=0$.
\end{theorem}

By Theorem~\ref{thm:Jchi} below each $J^T_\chi$ has
degree at most $1$, so $J^T(f)=\sum_\chi J^T_\chi(f)$, with
Theorem~\ref{thm:coarse}, is an identity of polynomials of degree at most
$1$ and holds for every $T$, in particular at $T=0$.

\begin{lemma}[Maass--Selberg relations]\label{lem:MS}
Let $\mu,\mu'\in\Pi_{\mathrm{aut}}(X_A)$, $\varphi\in\mathcal H_B(\mu)$,
$\varphi'\in\mathcal H_B(\mu')$, $T\ge0$, and let $\lambda,\lambda'$ be
regular points with $\lambda+\bar\lambda'\ne0$ and
$\lambda-\bar\lambda'\ne0$. Then, with $M=M(w,\lambda)$ and
$\check M=M(w,\lambda')$,
\begin{multline}\label{eq:MS}
\bigl\langle\Lambda^TE(\varphi,\lambda),\Lambda^TE(\varphi',\lambda')
\bigr\rangle
=\frac{\langle\varphi,\varphi'\rangle e^{(\lambda+\bar\lambda')T}
-\langle M\varphi,\check M\varphi'\rangle e^{-(\lambda+\bar\lambda')T}}
{\lambda+\bar\lambda'}\\
+\frac{\langle M\varphi,\varphi'\rangle e^{(-\lambda+\bar\lambda')T}}
{-\lambda+\bar\lambda'}
+\frac{\langle\varphi,\check M\varphi'\rangle e^{(\lambda-\bar\lambda')T}}
{\lambda-\bar\lambda'} .
\end{multline}
The inner products are those of~\eqref{eq:HBmu} and vanish unless the
characters agree; in particular the last two terms vanish unless
$\mu'=\mu^w$.
\end{lemma}

\begin{proof}[Sketch]
For $\mathrm{Re}\,\lambda'$ large, Lemma~\ref{lem:truncation}(b) gives
$\langle\Lambda^TE,\Lambda^TE'\rangle=\langle\Lambda^TE,E'\rangle$. Unfold
$E'$ and integrate over $N(\Q)\backslash N(\A)$ first. By the proof of
Lemma~\ref{lem:truncation}, $(\Lambda^TE)_B=E_B$ on $\{H\le T\}$ minus the
terms $\delta\ne1$; folding these back into a sum over $G(\Q)$ and
unfolding the other way, they contribute
$-\int_{H>T}E_B\,\overline{(\check M\varphi')_{-\lambda'}}$, because the
constant term of $E'-\varphi'_{\lambda'}$ is $(\check M\varphi')_{-\lambda'}$.
So
\[
\langle\Lambda^TE,E'\rangle=\int_{H(x)\le T}E_B\,\overline{\varphi'_{\lambda'}}
-\int_{H(x)>T}E_B\,\overline{(\check M\varphi')_{-\lambda'}},
\]
integrals over $N(\A)A(\Q)\backslash\GA$. Insert
Theorem~\ref{thm:eisenstein}(b); in the coordinates $x=\tilde a^1e^Xk$ each
of the four products is $e^{\nu X}$ times a product of values on
$\overline{K_0}$, with $\nu\in\{\pm\lambda\pm\bar\lambda'\}$, and
$\int_{-\infty}^Te^{\nu X}dX=e^{\nu T}/\nu$, $\int_T^\infty e^{\nu X}dX=
-e^{\nu T}/\nu$ in the appropriate half planes. The integral over
$X_A\times\overline{K_0}$ gives the inner products~\eqref{eq:HBmu}. The
general case follows by analytic continuation. In general the relations
are due to Langlands; see~\cite[\S4]{A3}, \cite{MW2}, \cite{LiIII}.
\end{proof}

\begin{theorem}\label{thm:Jchi}
Let $f\in\mathcal H(\GA)$, and for $\mu\in\Pi_{\mathrm{aut}}(X_A)$ and
$\lambda\in i\R$ put
\begin{equation}\label{eq:MT}
M_B^T(\mu,\lambda)=2T-M(w,\lambda)^{-1}\frac{d}{d\lambda}M(w,\lambda),
\end{equation}
an operator on $\mathcal H_B(\mu)$, self-adjoint for $\lambda\in i\R$.
\begin{enumerate}[label=(\alph*)]
\item If $\chi$ is the class of $(G,\pi)$, then $J^T_\chi(f)=
  m(\pi)\operatorname{tr}\pi(f)$.
\item If $\chi\in\mathcal X_A$ is unramified and $\mu\in\chi$, then
  \[
  J^T_\chi(f)=\frac1{2\pi}\int_\R\operatorname{tr}\bigl(M_B^T(\mu,it)\,
  I(\mu_{it},f)\bigr)\,dt .
  \]
\item If $\chi=\chi_b$ is ramified and $\mu=\chi_{\psi_b}$, then
  \[
  J^T_\chi(f)=\operatorname{tr}\bigl(R(f)|_{V_b}\bigr)
  +\frac12\cdot\frac1{2\pi}\int_\R\operatorname{tr}\bigl(M_B^T(\mu,it)\,
  I(\mu_{it},f)\bigr)\,dt
  +\frac14\operatorname{tr}\bigl(M(w,0)\,I(\mu_0,f)\bigr).
  \]
\end{enumerate}
\end{theorem}

\begin{proof}
By Theorem~\ref{thm:Jchi-poly}(c) it suffices to compute
$\int\Lambda^T_2K_{G,\chi}(x,x)\,dx$ up to $o(1)$ as $T\to\infty$. By
Lemma~\ref{lem:truncation}(b) and Proposition~\ref{prop:kernels-spec}
this is $\sum_\phi\langle R(f)\phi,\Lambda^T\phi\rangle$ in case~(a) and
\[
\sum_\theta\langle R(f)\theta,\Lambda^T\theta\rangle+\frac12\sum_{\mu\in\chi}
\frac1{2\pi}\int_\R\sum_{\varphi\in\mathcal B(\mu)}\bigl\langle\Lambda^T
E(I(\mu_{it},f)\varphi,it),\Lambda^TE(\varphi,it)\bigr\rangle\,dt
\]
in cases~(b)--(c), the inner product on $i\R$ being the value at
$\lambda'=\lambda$ of~\eqref{eq:MS}, which is regular there; the
interchange of the integrals is justified as in~\cite[\S4]{A3}. (a) Cusp
forms satisfy $\Lambda^T\phi=\phi$. For the residual part,
$\Lambda^T\theta\to\theta$ in $L^2$, so the first sum tends to
$\operatorname{tr}(R(f)|_{V_b})$.

Now let $\lambda=it$, $\lambda'=\lambda+\varepsilon$ with
$\varepsilon\in i\R$, so $\lambda+\bar\lambda'=-\varepsilon$. Expanding
$M(w,\lambda')=M+\varepsilon\,\frac{dM}{d\lambda}+O(\varepsilon^2)$ and
using $M^*=M^{-1}$ on $i\R$, the first fraction of~\eqref{eq:MS} tends to
$\langle\varphi,M_B^T(\mu,\lambda)\varphi'\rangle$ as $\varepsilon\to0$.
With $\varphi=I(\mu_\lambda,f)\varphi_j$, $\varphi'=\varphi_j$ and the sum
over $j$ this gives $\operatorname{tr}(M_B^T(\mu,\lambda)
I(\mu_\lambda,f))$, using that $M^{-1}\frac{dM}{d\lambda}$ is self-adjoint on
$i\R$ (differentiate $M(it)^*M(it)=1$). For $\mu^w\ne\mu$ the two $\mu\in
\chi$ give equal contributions, and the factor $\frac12$ cancels; this
proves~(b), the expression being already linear in $T$.

For $\mu^w=\mu$ the last two terms of~\eqref{eq:MS} also contribute, with
$-\lambda+\bar\lambda'\to-2it$ and $\lambda-\bar\lambda'\to2it$. Put
$F(t)=\operatorname{tr}(M(w,it)I(\mu_{it},f))$; by
Theorem~\ref{thm:eisenstein}(c),(e),
$\operatorname{tr}(M(w,-it)I(\mu_{it},f))=F(-t)$. The contribution is
\[
\frac12\cdot\frac1{2\pi}\,\mathrm{PV}\!\int_\R\frac{F(-t)e^{2itT}-F(t)
e^{-2itT}}{2it}\,dt
=\frac1{4\pi}\Bigl(\int_\R F(t)\frac{\sin2tT}{t}\,dt
+i\,\mathrm{PV}\!\int_\R F(t)\frac{\cos2tT}{t}\,dt\Bigr),
\]
which tends to $\frac1{4\pi}\cdot\pi F(0)=\frac14F(0)$ as $T\to\infty$, by
Dirichlet's integral and the Riemann--Lebesgue lemma applied to
$(F(t)-F(-t))/t$. Since $J^T_\chi$ is a polynomial, it equals the
polynomial part of these limits, which is~(c).
\end{proof}

Two remarks on the constants. First, $M_B^T$ is the case of
$SL(2)$ of Arthur's formula
\[
M_B^T=\mathrm{vol}(\mathfrak a/\Z\alpha^\vee)\lim_{\zeta\to0}
\sum_{s\in\{1,w\}}\frac{e^{(s\zeta)(T)}M(s,\lambda)^{-1}
M(s,\lambda+\zeta)}{\theta(s\zeta)},
\]
in which $\theta(\zeta)=\zeta(\alpha^\vee)=d\zeta$ and
$\mathrm{vol}(\mathfrak a/\Z\alpha^\vee)=d$ for our $dX$
(Remark~\ref{rem:alpha}); the factor $d$ cancels and the limit
is~\eqref{eq:MT} for every $k$. Proposition~\ref{prop:Tcoef} confirms the
coefficient of $T$ against the geometric side. Second, in Arthur's fine
expansion~\cite[Thm.~8.2]{A14}, \cite{LiIII} the last term of~(c) is
$|W^A_0||W^G_0|^{-1}|\det(w-1)_{\mathfrak a}|^{-1}
\operatorname{tr}(M(w,0)I(\mu_0,f))$, with $|W^A_0||W^G_0|^{-1}=\frac12$
and $|\det(w-1)_{\mathfrak a}|=2$. In the trace formula for $GL(2)$
of~\cite[\S6]{GJ} the same term is written $-\frac14\operatorname{tr}
(M(0)\pi_0(f))$, next to $+\frac1{4\pi}\int\operatorname{tr}(M(-iy)M'(iy)
\pi_{iy}(f))\,dy$: that formula has the trace of $R(f)$ on the discrete
spectrum alone on one side, and moving both terms to the spectral side
gives the signs of Theorem~\ref{thm:tf}(e). The sign of the cross terms
in~\eqref{eq:MS}, from which it comes, is also checked numerically on
the Eisenstein series of $SL(2,\Z)$ (Remark~\ref{rem:verified-sec6}),
and $M(w,0)$ is computed in Proposition~\ref{prop:Mw0}.

\begin{corollary}\label{cor:mprime}
Let $S$ and $\mu=\chi_\psi\cdot(\chi\circ p)$ be as in
Lemma~\ref{lem:intertwining}, $f=f_S\otimes e^S$, and $\lambda\in i\R$.
Then $I(\mu_\lambda,f)=0$ unless $\mu_v$ is trivial on $A(\OO_v)^*$ for
all $v\notin S$, and in that case
\begin{multline*}
\operatorname{tr}\Bigl(M(w,\lambda)^{-1}\tfrac{dM}{d\lambda}(w,\lambda)
I(\mu_\lambda,f)\Bigr)\\
=\frac{(m^S)'}{m^S}(\mu,\lambda)\operatorname{tr}
I_S(\mu_{S,\lambda},f_S)+\operatorname{tr}\Bigl(M_S(\lambda)^{-1}
\tfrac{dM_S}{d\lambda}(\lambda)I_S(\mu_{S,\lambda},f_S)\Bigr),
\end{multline*}
where $I_S=\bigotimes_{v\in S}I_v$ and
$\frac{(m^S)'}{m^S}(\mu,\lambda)=2\frac{(L^S)'}{L^S}(2\lambda,\chi^2)
-2\frac{(L^S)'}{L^S}(2\lambda+1,\chi^2)$.
\end{corollary}

\begin{proof}
$I(\mu_\lambda,e_v)$ is the projection onto the $K_v^*$-fixed vectors,
which are $0$ unless $\mu_v$ is trivial on $A(\OO_v)^*$. Then apply
Lemma~\ref{lem:intertwining}(a)--(b) and the Leibniz rule.
\end{proof}

\begin{example}\label{ex:mprime-Q}
For $k=\Q$, $S=\{\infty,2\}$ and $\mu=\chi_\psi$,
$\frac{(m^S)'}{m^S}(\lambda)=2\frac{\zeta'}{\zeta}(2\lambda)
-2\frac{\zeta'}{\zeta}(2\lambda+1)+\frac{2\log2\cdot2^{-2\lambda}}
{1-2^{-2\lambda}}-\frac{2\log2\cdot2^{-1-2\lambda}}{1-2^{-1-2\lambda}}$.
For $SL(2)$ the corresponding factor is
$\frac{\zeta'}{\zeta}(\lambda)-\frac{\zeta'}{\zeta}(\lambda+1)$ plus the
$2$-adic terms.
\end{example}

\subsection{The trace formula}\label{sec:tf}

\begin{theorem}\label{thm:tf}
Let $S$ be a finite set of places of $\Q$ as in
Section~\ref{sec:local-dist}, large enough for Theorem~\ref{thm:fine} and
containing the places below which $\psi$ is ramified; let $f_S\in
\mathcal H(\overline{G(\Q_S)})$ and $f=f_S\otimes e^S$. Then
$\sum_{\mathcal O}J_{\mathcal O}(f)=\sum_{\chi}J_\chi(f)$, where both
sides converge absolutely and:
\begin{enumerate}[label=(\alph*)]
\item for the central classes $\mathcal O_z$, $z\in\{\pm1\}$,
  \[
  J_{\mathcal O_z}(f)=a^G(S,1)f_S(z)+\sum_{\xi\in k_S^\times/k_S^{\times2}}
  a^G(S,u_\xi)\,J_G(zu_\xi,f_S)+\tfrac12\,c_k\,J_A(z,f_S),
  \]
  with the coefficients $a^G(S,u_\xi)$ of Theorem~\ref{thm:fine}, which
  are computed in Proposition~\ref{prop:unipotent};
\item for $\mathcal O$ elliptic, $J_{\mathcal O}(f)=
  \mathrm{vol}(T_\mu(\Q)\backslash T_\mu(\A))\,J_G(\mu_S,f_S)$;
\item for $\mathcal O$ hyperbolic, the class of $a(t)$,
  $J_{\mathcal O}(f)=c_k\,J_A(a(t)_S,f_S)$;
\item for $\chi$ the class of $(G,\pi)$, $J_\chi(f)=m(\pi)
  \operatorname{tr}\pi(f)$;
\item summed over $\chi\in\mathcal X_A$,
  \begin{multline*}
  \sum_{\chi\in\mathcal X_A}J_\chi(f)=\sum_{b}\operatorname{tr}
  \bigl(R(f)|_{V_b}\bigr)
  -\frac12\sum_{\mu}\frac1{2\pi}\int_\R\operatorname{tr}\Bigl(M(w,it)^{-1}
  \tfrac{dM}{d\lambda}(w,it)\,I(\mu_{it},f)\Bigr)\,dt\\
  +\frac14\sum_b\operatorname{tr}\bigl(M(w,0)\,I(\chi_{\psi_b,0},f)
  \bigr),
  \end{multline*}
  where $b$ runs over $k^\times/k^{\times2}$ and $\mu$ over
  $\Pi_{\mathrm{aut}}(X_A)$; only finitely many $b$, and only $\mu$
  unramified outside $S$, contribute (Example~\ref{ex:ramified-Q},
  Corollary~\ref{cor:mprime}).
\end{enumerate}
\end{theorem}

\begin{proof}
At $T=0$, Theorem~\ref{thm:coarse}, Theorem~\ref{thm:fine} and
Theorem~\ref{thm:Jchi-poly}(b) give $J(f)=\sum_{\mathcal O}J_{\mathcal O}(f)
=\sum_\chi J_\chi(f)$. (a)--(c) are Theorem~\ref{thm:fine}. (d)--(e) are
Theorem~\ref{thm:Jchi} at $T=0$: in~(b) of that theorem, write the
integral as half the sum over the two members of $\chi$; in~(c) the class
has one member. The absolute convergence of the spectral side
is~\cite{A3}, \cite{LiIII}.
\end{proof}

\begin{proposition}\label{prop:Tcoef}
The coefficients of $T$ in $\sum_\chi J^T_\chi(f)$ and in
$\sum_{\mathcal O}J^T_{\mathcal O}(f)$ agree. Both equal
$c_k\sum_{\gamma\in A(\Q)}F_f(\gamma)$, where $F_f(x)=\int_{\overline{K_0}}
\int_{N(\A)}f(\tilde k^{-1}xn\tilde k)\,dn\,d\tilde k$.
\end{proposition}

\begin{proof}
Spectral side. By Theorem~\ref{thm:Jchi}, only the term $2T$ of $M_B^T$
contributes, and summing over $\chi$ the coefficient is
$\sum_\mu\frac1{2\pi}\int_\R\operatorname{tr}I(\mu_{it},f)\,dt$ (each
unramified class contributes twice $\frac12$ through its two members).
Computing $I(\mu_\lambda,f)$ on $\overline{K_0}$ in the coordinates
$y=n\tilde ak'$ of~\eqref{eq:iwasawa} gives the kernel
$\int_{\overline{A(\A)}}\int_{N(\A)}f(\tilde k^{-1}\tilde an\tilde k')\mu(\tilde
a)e^{(\lambda+1)H(\tilde a)}\,dn\,d\tilde a$, so $\operatorname{tr}
I(\mu_\lambda,f)=\int_{\overline{A(\A)}}e^{H(\tilde a)}F_f(\tilde a)
\mu(\tilde a)e^{\lambda H(\tilde a)}\,d\tilde a$. Fourier inversion on $\R$
($dX$, $dt/2\pi$) and then Theorem~\ref{thm:tf-rank0} on $X_A$, applied
to the anti-genuine function $F_f$ on $\overline{A(\A)^1}$, give
$\sum_\mu\frac1{2\pi}\int\operatorname{tr}I(\mu_{it},f)\,dt=c_k\sum_{\gamma\in
A(\Q)}F_f(\gamma)$.

Geometric side. By Propositions~\ref{prop:hyperbolic}--\ref{prop:central},
the coefficient is $c_k\bigl(\hat\Phi_f(0)+\hat\Phi_{f_{-1}}(0)\bigr)
+2c_k\sum_{\mathcal O\ \mathrm{hyp}}\int_{\overline{A(\A)}\backslash\GA}
f(y^{-1}\mu y)\,dy$. Here $\hat\Phi_{f_z}(0)=F_f(z)$ for $z=\pm1$. For
$\mu=a(t)$, $t\in k^\times\smallsetminus\{\pm1\}$, the coordinates
$y=n(z)\tilde k$ and $y^{-1}\mu y=\tilde k^{-1}\mu\,n((1-t^{-2})z)\tilde k$
give $\int f(y^{-1}\mu y)\,dy=|1-t^{-2}|_{\A_k}^{-1}F_f(\mu)=F_f(\mu)$ by the
product formula. The integral is the same for $a(t^{-1})=w\,a(t)\,w^{-1}$
(Lemma~\ref{lem:w-action}(a)), so $2F_f(a(t))=F_f(a(t))+F_f(a(t^{-1}))$,
and the hyperbolic class of $a(t)$ accounts for the two elements
$a(t^{\pm1})$ of $A(\Q)$.
\end{proof}

Proposition~\ref{prop:Tcoef} is a consistency check between
Sections~\ref{sec:geometric} and~\ref{sec:spectral}: with $4T$ in place
of $2T$ in $M_B^T$, or with $n(A_B)=1$ in place of $2$, the spectral
coefficient would be twice the geometric one.

\begin{remarkrev}\label{rem:sec5-norm}
\emph{Normalizations.} Everything here uses $\lambda(X)=\lambda X$ on
$\mathfrak a\cong\R$, Lebesgue $dX$ and $dt/2\pi$ on $i\R$, and the inner
product~\eqref{eq:HBmu}, which includes $c_k$. Arthur normalizes the
measure on $\mathfrak a$ so that $\mathfrak a/\Z\alpha^\vee$ has volume
$1$, i.e.\ $dX/d$, and takes the dual measure on $i\mathfrak a^*$; the
operator $M_B^T$ then carries the factor $\mathrm{vol}(\mathfrak a/
\Z\alpha^\vee)\theta(\zeta)^{-1}$, and the products in
Theorem~\ref{thm:Jchi} do not change (Remark~\ref{rem:alpha}). In the
classical variable $s=(\lambda+1)/2$ of the Conventions above, for
$k=\Q$, $\frac{d}{d\lambda}=\frac12\frac{d}{ds}$, and $2T=\log y_T$, where
$y_T=e^{2T}$ is the height at which $\Lambda^T$ truncates.
\end{remarkrev}

\subsection{Status of the results}\label{sec:status5}

\begin{remarkrev}\label{rem:sec5-status}
\emph{What rests on citation.} Proved here: Lemma~\ref{lem:factor} (given Flath's
theorem); Lemmas~\ref{lem:w-action}, \ref{lem:truncation}(a)--(b)
and~\ref{lem:kchi}; Lemma~\ref{lem:intertwining}, with~(b) also checked
numerically;
Theorem~\ref{thm:residual}(a)--(b), the computations in the proofs of
Lemma~\ref{lem:MS} and Theorem~\ref{thm:Jchi} from~\eqref{eq:MS}, and
Proposition~\ref{prop:Tcoef}. Taken from the literature:
Lemma~\ref{lem:spec-decomp}, Proposition~\ref{prop:Lchi},
Theorem~\ref{thm:eisenstein}~\cite{L2,MW2}; the wave-packet formula~\eqref{eq:KGchi} and the
absolute convergence in Proposition~\ref{prop:kernels-spec};
Lemma~\ref{lem:truncation}(c); Theorem~\ref{thm:Jchi-poly}~\cite{A3,
LiIII}, including the interchange of integrals used for
Theorem~\ref{thm:Jchi}; the analytic continuation in Lemma~\ref{lem:MS};
and Theorem~\ref{thm:residual}(c)~\cite{Ge2,GPS80,Wa80}, except for
the matching $V_b\cong r^+_{\psi_{-b}}$ of the parameters, which is proved
here from Lemma~\ref{lem:weil-rep}. The matching depends on the
normalization~\eqref{eq:weil-rep}: if $n(u)$ acts by $\psi'(\frac12ux^2)$,
it reads $V_b\cong r^+_{\psi_{-2b}}$. We use from the cited papers that
the residual spectrum is the sum of the spaces of even theta
functions~\eqref{eq:theta} and that these are irreducible. The range ``$T$ sufficiently large'' in
Theorem~\ref{thm:Jchi-poly} is Arthur's; for this group we expect $T\ge0$
to suffice, as on the geometric side, but have not checked it.
\end{remarkrev}

\begin{remarkrev}\label{rem:verified-sec5}
The script \texttt{checks/sec5\_n2\_check.py} verifies, for $k=\Q$:
Lemma~\ref{lem:w-action}(a)--(c) at $v=\infty,2,3,5,7$ (exactly, with
$\sigma_v$ and the Weil index); the unramified intertwining integral of
Lemma~\ref{lem:intertwining}(b) at $p=3,5,7$, by an explicit
decomposition of the lift of $w^{-1}n(y)$, shell by shell, including the
shape of the integrand on odd and even shells, and the same integral for
$SL(2)$; Arthur's limit for $M_B^T$ for $d=1,2,3$ and the 1995
normalization (symbolically); the diagonal limit of~\eqref{eq:MS} and the
self-adjointness of $M^{-1}M'$ for a model unitary family; the limit
$\frac14F(0)$ of the cross terms; the residue $2/\pi^2$ of
Example~\ref{ex:theta}; at $v=\R$, the two local identities behind
Proposition~\ref{prop:Tcoef}; and, for Lemma~\ref{lem:weil-rep}, the
product $(w(t),1)(w(1),1)^{-1}=(a(t),(-1,t)_v)$ of unipotent lifts and the
identity $(\gamma_v(t)/\gamma_v(1))^2=(t,-1)_v$ at $v=\infty,2,3,5,7$, and
the weights $\pm\frac12$ of the ground state at $v=\R$.
\end{remarkrev}
